% The ten-step coordination storyline as message flow between the
% organisations, plus the conditional stand-down. Current design (ADR-0009 to
% ADR-0013), not the pre-ADR-0011 bus-driver sequence diagram.
% Sources: thesis-pack/04-architecture-agentic/agents/bis.md (action table),
% adr/0011 (transport_plan), adr/0012 (facility cleared), adr/0013 (shelter
% edge), adr/0009 (stand-down). Input from inc/casestudy.tex.
\begin{figure}[htbp]
\centering
\begin{tikzpicture}[x=1.0cm, y=1.0cm,
  llm/.style={draw, rounded corners=2pt, fill=blue!8, align=center,
    font=\scriptsize, text width=1.6cm, minimum height=0.9cm, inner sep=2pt},
  scr/.style={draw, rounded corners=2pt, fill=gray!12, align=center,
    font=\scriptsize, text width=1.6cm, minimum height=0.9cm, inner sep=2pt},
  life/.style={densely dashed, gray!60},
  msg/.style={-{Stealth[length=1.6mm]}, thick},
  alt/.style={-{Stealth[length=1.6mm]}, dashed, gray!70!black},
  lbl/.style={above, font=\scriptsize, fill=white, inner sep=1pt, yshift=-1pt},
  note/.style={draw, rounded corners=2pt, fill=orange!8, align=center,
    font=\scriptsize, text width=1.7cm, inner sep=2pt},
]
  % Lanes: x position, style, header text.
  \foreach \x/\s/\t in {
      0/scr/{\textbf{ZKD}\\command},
      2/llm/{\textbf{BIS}\\coordinator},
      4/scr/{\textbf{Augustinum}\\nursing home},
      6/scr/{\textbf{SC1, SC2}\\shelters},
      8/llm/{\textbf{Hochbahn}\\transport},
      10/llm/{\textbf{Bus drivers}\\bus-1, bus-2},
      12/llm/{\textbf{Polizei}\\escort},
      14/llm/{\textbf{DRK}\\ambulance}} {
    \node[\s] at (\x,0) {\t};
    \draw[life] (\x,-0.45) -- (\x,-11.7);
  }

  % 1 Alarm
  \draw[msg] (0,-1.0) -- node[lbl] {1 alarm} (2,-1.0);
  % 2 Assessment
  \draw[msg] (2,-1.6) -- node[lbl] {2 residents?} (4,-1.6);
  \draw[msg] (4,-2.1) -- node[lbl] {100, deadline} (2,-2.1);
  % 3 Resource allocation
  \draw[msg] (2,-2.7) -- node[lbl, pos=0.3] {3 free beds?} (6,-2.7);
  \draw[msg] (6,-3.2) -- node[lbl, pos=0.7] {60 + 60} (2,-3.2);
  % 4 Route planning: done by the drivers, not by a coordinator
  \node[note] at (10,-3.95) {4 route legs on real roads};
  % 5 Logistics
  \draw[msg] (2,-4.7) -- node[lbl, pos=0.35] {5 request buses} (8,-4.7);
  \draw[msg] (8,-5.2) -- node[lbl, pos=0.65] {transport plan} (2,-5.2);
  \draw[msg] (8,-5.7) -- node[lbl] {runs} (10,-5.7);
  % 6 Police escort, 7 medical support
  \draw[msg] (2,-6.3) -- node[lbl, pos=0.85] {6 escort} (12,-6.3);
  \draw[msg] (2,-6.8) -- node[lbl, pos=0.9] {7 medical} (14,-6.8);
  % 8 Execution, repeated per run
  \draw[msg] (10,-7.4) -- node[lbl] {8 departed, $n$ aboard} (4,-7.4);
  \draw[msg] (4,-7.9) -- node[lbl] {facility cleared} (2,-7.9);
  % 9 Arrival and handover, repeated per run
  \draw[msg] (10,-8.5) -- node[lbl] {9 drop-off} (6,-8.5);
  \draw[msg] (10,-9.0) -- node[lbl] {drop-off} (8,-9.0);
  \draw[msg] (8,-9.5) -- node[lbl, pos=0.6] {transport complete} (2,-9.5);
  % 10 Mission report
  \draw[msg] (2,-10.1) -- (0,-10.1);
  \node[font=\scriptsize, anchor=west] at (2.1,-10.1) {10 mission complete};
  % Conditional stand-down (ADR-0009), exercised by no scenario
  \draw[alt] (2,-10.8) -- (0,-10.8);
  \node[font=\scriptsize, anchor=west, text=gray!50!black] at (2.1,-10.8)
    {stand-down only: request abort};
  \draw[alt] (0,-11.3) -- (2,-11.3);
  \node[font=\scriptsize, anchor=west, text=gray!50!black] at (2.1,-11.3)
    {countersigned, then recall of transport and escorts};

  % Legend
  \node[llm, text width=0.6cm, minimum height=0.35cm] at (9.6,-12.15) {};
  \node[font=\scriptsize, anchor=west] at (10.0,-12.15) {LLM-driven};
  \node[scr, text width=0.6cm, minimum height=0.35cm] at (12.0,-12.15) {};
  \node[font=\scriptsize, anchor=west] at (12.4,-12.15) {scripted};
\end{tikzpicture}
\caption[The coordination storyline]{The coordination storyline as message
flow between the organisations of the case. Numbers give the storyline step.
Steps 8 and 9 repeat for every bus run. Step 4 is not a message: each driver
routes its own legs over real roads. The dashed branch is the conditional
stand-down, which no evaluated scenario triggered. In the workflow baseline, one
Incident Commander takes the place of the two coordinator lanes, BIS and
Hochbahn. All other lanes are shared by both systems.}
\label{fig:cs:storyline}
\end{figure}
